% Transcription — fonds Alexandre Grothendieck, shelfmark 37, batch 1 (pages 1-20).
% Established from the facsimile archives/batches/37/batch-01.pdf (a copy of
% 37.pdf fetched through the project's relay, grothendieck-relay.onrender.com;
% PDF page k = archivists' page k, no cover sheet), per the skill
% transcribe-grothendieck. The folder is 99 pages in five batches (1-20, 21-40,
% 41-60, 61-80, 81-99); this is the first, done in parallel with the others.
%
% Pass: Opus 5.5 (claude-opus-5-5), 2026-09-26 — first pass, unchecked against the pages by a human.
%
% Pages skipped: 1 (a blank cream sheet). Pages 5 and 7 are tan covers
% inscribed in pencil, top right, « SGA 1 » and « SGA 4 »: no \page marker,
% their words become section headings with a \note (hand.md §5; the same
% treatment as the « SGA 5 », « SGA 6 » covers of batch 3).
% Pages noted only (not re-set): 8, 9, 10, 11, 13.
% Pages transcribed: 2, 3, 4, 6, 12, 14, 15, 16, 17, 18, 19, 20.
%
% Decisions, item by item, with the evidence:
% - Page 2: a typed sheet « SGA : Manuscripts à frapper » — a list, seminar by
%   seminar (SGA 1, 4, 5, 6, 7), of the exposés still to be typed, with their
%   rédacteurs and page estimates, and an NB on the typing and the
%   secretary's fee. DECISION: his own unpublished typescript, transcribed in
%   full. Evidence: it is written in the first person by the man organising
%   the re-editions (« Je me mets en rapport avec lui [Verdier] ces
%   jours-ci »; « cela dépendra de Mr. Verdier, qui est à Brandeis cette
%   année », matching his letters of 3.11.1970 on pp. 3, 4, 15 and Verdier's
%   Brandeis letterhead on p. 13); it lists « Exp XIV Grothendieck » as one
%   item among others; it is an instruction sheet, not a text of any seminar.
%   The pencil marks (the SGA 1 line struck, « écrire Bucur », « écr.
%   Deligne », a box, brackets, a wavy stroke cancelling the SGA 6 block,
%   ticks) are taken as his — the same pencil hand as p. 17 and p. 8 —
%   presumably, not provably. The pencil name after the SGA 1 line is not
%   read (last word perhaps « Saarbach »).
% - Pages 3, 4: carbons of his letters, typed, « Massy le 3.11.1970 », to
%   Madame Michon (missing IHES archive copies of SGA 4 XVII and VI, needed
%   for the photo-offset edition) and to Deligne (Rim's SGA 7 VI, Mlle
%   Altazin's typing of SGA 4 XVIII, Ferrand's SGA 6 XI). Transcribed in full.
% - Page 6: a typed leaf paginated « (ix) », signed off « Massy, Août 1970 »:
%   the last page of a preface by him (military subsidies of the IHES, his
%   leaving it, Springer and K. Peters). DECISION: transcribed in full as an
%   unpublished state of his own text. Evidence: compared on 2026-09-26 with
%   the published Préface of SGA 1 (LNM 224; arXiv math/0206203, pp. xi-xii):
%   the published text has « immoral et dangereux » with a footnote (4), not
%   « profondément immoral et extrêmement dangereux »; « quitte l'IHES le 30
%   septembre 1970 », not « le 30 Septembre prochain »; « J'ai demandé à
%   M. Motchane ... que l'IHES s'abstienne à partir du 1er octobre 1970 de
%   diffuser », not « Monsieur L. Motchane ... s'est déclaré d'accord pour
%   qu'à partir du 1er Octobre 1970, l'IHES ne diffuse plus »; and it adds a
%   paragraph thanking J.P. Delale for the indexes, absent here. Nothing in
%   his hand on the leaf.
% - Page 8: a typed letter of J.-L. Verdier (signed with his paraph, the
%   same as on p. 13), dated by hand « 16 Juillet » (no year), sending « la
%   nouvelle rédaction de 4.10.6 », a list of misprints in SGA 4 IV nos 9
%   and 10 sqq., and the news that no 14 is missing. Another person's letter:
%   one \note; his pencil « demander à Messing » in the margin is given in
%   full.
% - Pages 9-11: the typescript that letter encloses — Exercices 4.10.6 and
%   4.10.7 of SGA 4 IV (petit site S(X), Res_X, Prol_X, the embedding of the
%   petit topos in the gros topos), typed pages with the machine's symbols
%   inked in. Verdier's redaction, not his text: one \note per page, not
%   re-set. The only handwritten marks are proof-reading ones (an insertion
%   sign for the typed marginal « le site » on p. 9; an ink « (V 4.1) » called
%   in on p. 10, a handwritten S(X)~ at a line end); whose they are is not
%   established and they are given in full. Page 11 carries none.
% - Page 12: a photocopy, clipped at both edges, of his letter « Massy le
%   23.6.1971 » to Verdier (fascicule 1 of SGA 4 held up by exercise 4.10.6).
%   Transcribed in full; the letters lost at the edges are \supplied{}.
% - Page 13: a handwritten letter of Verdier on Brandeis University
%   letterhead, « le 28 Mai » (no year), answering p. 14 (the two appendices,
%   a « version non canulée » of 4.10.6, IV §§9-13, a P.S. on Mike [Artin]).
%   Another person's letter: one \note; his only mark, a faint pencil oval
%   round « de 4.10.6 », is given (attribution presumable).
% - Pages 14, 15: his typed letters to Verdier, « Massy le 18.5.1971 » and
%   « Massy le 3.11.1970 ». Transcribed in full. Page 14 continues on a leaf
%   not in this batch (a closing line shows through from the back).
% - Page 16: his manuscript on IHES letterhead, « Continuation de SGA 4 IV
%   4.10.6 »: a 5°) for the exercise (a full subcategory S_0 of S, conditions
%   (i), (ii), M_0, Res and Prol). Transcribed in full; fast hand, many \ill.
% - Page 17: a typed « SGA 4 VI / Sommaire » (ten numbers) followed by an NB
%   and two paragraphs in the first person on what nos 4 and 8 will contain
%   and on the state of the redaction of nos 1-3, heavily revised in his ink
%   (renumbering, « 4. Constructibilité » replaced, « 10. Appendice (par
%   P. Deligne) », page ranges 1 à 19 ... 47 à 65, an added example) and in
%   pencil (the NB cancelled by diagonals, a note at the head). DECISION:
%   transcribed in full as his own unpublished typescript. Evidence that it
%   is not the published text: the published Exposé VI (LNM 270; the
%   Orgogozo re-edition's table, consulted 2026-09-26) has an Introduction,
%   « 5. Commutation des foncteurs H^i(X,-) aux limites inductives
%   filtrantes », « 6. Limites inductive et projective d'une catégorie
%   fibrée », no « Constructibilité », no NB, no statement of intentions.
%   Evidence that it is his: the author announces he will write no 4
%   himself (« Je pense que je le rédigerai mais ne le garantis pas »),
%   while his letter of 3.11.1970 (p. 15) assigns VI « par. 5 et suivants »
%   to Verdier; the ink revisions are in the hand of pp. 18-19. Presumable,
%   not proven.
% - Pages 18-20: his own manuscript notes, transcribed in full: p. 18 a list
%   of reference checks in SGA 4 XVII (pp. 75-175 of a typescript); p. 19 a
%   plan of SGA 4 in three fascicules with page counts per exposé (totals
%   432, 178, 640, « # 570 p. »); p. 20 instructions on a missing passage to
%   be inserted as a « p. 83 bis » of a stencil, corrections Lg^* / Lg'^* on
%   p. 85, and an addition 210 + 76 + 250 = 536; the upper part is cancelled
%   by three pencil diagonals and « Le texte oublié a été reporté ici par
%   erreur !!! ».
%
% His pagination and numbered runs: none of his own on his manuscript pages.
% Typed « (ix) » on p. 6. On p. 17 the sommaire's numbers 1-10 (typed 1-9,
% renumbered by hand) and his page ranges « 1 à 19, 20 à 37, 38 à 46, 47 à
% 65 ». Page 16 numbers its item « 5° », continuing 1°-4° of exercise
% 4.10.6. Dates on the leaves: Août 1970 (p. 6); 3.11.1970 (pp. 3, 4, 15);
% 18.5.1971 (p. 14); « 28 Mai » without year (p. 13); 23.6.1971 (p. 12,
% which thanks for « ta lettre du 23 Mai »); « 16 Juillet » without year
% (p. 8).
%
% What the batch is about. The logistics of the Springer re-edition of the
% SGA in 1970-1971, after his departure from the IHES: what remains to be
% typed and by whom (p. 2), the preface's last page in a state earlier than
% the printed one (p. 6), the chase for the IHES archive copies, for
% Deligne's appendices and for a corrected version of the « canulé »
% exercise SGA 4 IV 4.10.6 (pp. 3-4, 8-16), the plan of SGA 4 VI (p. 17),
% reference checks in SGA 4 XVII (p. 18), the division of SGA 4 into three
% fascicules (p. 19) and a correction to a stencil (p. 20). The mathematics
% is on p. 16 (a generalisation of 4.10.6 3°-4° to a full subcategory S_0
% of S stable under M) and in the sommaire of p. 17.
%
% Notation fixed once: \underline{M} as the typescripts underline it;
% S^{\sim} for the topos of sheaves (his tilde); \operatorname{ob},
% \mathrm{Fl}; the arrowed Lim of p. 17 as \varinjlim. Readings to check
% first: the pencil name on p. 2; the struck words and the line « Munissons
% ... cocontinu » on p. 16; the struck handwritten line and its interlinear
% replacement on p. 17 (item 10); « V ... Deligne » on p. 18; the fascicule
% 2 and 3 titles on p. 19; « me montrer la frappe faite » on p. 20.

\documentclass[11pt,a4paper]{article}
\input{../preamble/grothendieck}

\folder{37}
\batch{1}
\pages{1}{20}
\dating{1967-1971}
\watermark{Édition de démonstration}
\foldertitle{Rééditions SGA, EGA [SGA 4 à SGA 7, EGA I et EGA II] : tapuscrits et copies de tapuscrit annoté (s.d.), lettres (1967-1971, s.d.), notes manuscrites (s.d.)}

\begin{document}

\page{2}
\note{feuille dactylographiée, sans date ni signature, qu'on tient pour un
tapuscrit de lui, inédit : il y parle à la première personne en organisateur
des rééditions (« Je me mets en rapport avec lui ces jours-ci », à propos de
Verdier, « qui est à Brandeis cette année », comme dans ses lettres du
3.11.1970, pages 3, 4 et 15). Les marques au crayon sont données comme ajouts
(entre coins) ou en marge ; elles sont vraisemblablement de sa main, sans que
cela soit établi. Les surfrappes de la machine ne sont pas relevées}

SGA : Manuscripts à frapper

\begin{itemize}
\item[SGA 1] \struck{Index notations et terminologique : \emph{10} p}
\marginal{\ill{} \uncertain{Saarbach}}
\item[SGA 4] Exp IV (fin), Verdier : \emph{50}. --- Exp V, Verdier :
\emph{80}. --- Exp VI (fin), Verdier : 50. --- Exp VI B, Saint Donat : 10. ---
Tables des matières pour les trois volumes, index : 30. --- Appendice à XVII,
Saint Donat : 10.
\item[SGA 5] Exp IV, ? : 100 ?. --- Exp VII, Jouanolou : \emph{50}. --- Exp XI,
Bucur : \emph{70}. --- Exp XII, Bucur : 100. --- Exp XIV, Grothendieck : 20.
--- Introduction, table des matières, index : 15.
\marginal{\add{écrire Bucur}}
\item[SGA 6] \struck{Exp XI, Ferrand : 70 ?. --- Table des matières, index : 10.}
\item[SGA 7] Exp I à XIV, ? : 120 ?. --- Exp V, Mme Raynaud : 30. --- Exp VI,
Rim : \emph{70}. --- Introduction, table des matières, index : 15.
\marginal{\add{écr. Deligne}}
\end{itemize}
\note{mise en colonnes du tapuscrit (rubrique, exposé, rédacteur, nombre de
pages) ramenée ici à une liste ; les nombres en italique sont soulignés à la
machine. Au crayon : un trait biffe toute la ligne de SGA 1, et deux mots
suivent le tiret en marge ; une accolade réunit les lignes « Exp V » à
« Appendice à XVII » de SGA 4, une plus petite « Exp VI (fin) » et
« Exp VI B » ; un crochet réunit « Exp XII » et « Exp XIV » de SGA 5, avec en
marge « écrire Bucur » ; un trait ondulé barre le bloc de SGA 6 ; un cadre
entoure les nombres 30 et 70 de SGA 7 (Exp V et VI), avec en marge « écr.
Deligne » souligné ; des traits obliques pointent en marge les lignes
« Exp V » (SGA 4), « Exp VII », « Exp XI », « Exp XII » (SGA 5) et
« Introduction » (SGA 7)}

NB Les chiffres soulignés correspondent à des exposés dont le manuscrit est
disponible dès maintenant ; ceux suivis d'un signe ? à ceux dont il n'est pas
certain qu'ils seront rédigés. En plus de la frappe d'exposés, tables de
matières, index, il y a pour chacun des cinq séminaires envisagés un certain
travail de mise au point (correction de fautes de frappe et autres erreurs,
insertion de notes de bas de page etc), représentant un travail variable d'un
séminaire à l'autre. Il faudrait sans doute fixer un forfait pour chacun d'eux,
pour la rénumération de la sécrétaire, après examen avec elle des corrections à
faire \ldots

Il y a actuellement dans les 250 pages prêtes pour la frappe. J'espère que dans
les deux ou trois mois qui viennent, il y aura aussi les env. 200 pages qui
manquent encore pour SGA 4 (qui sont déjà partiellement rédigés \ldots), ce qui
ferait dans les 450 à 500 pages à frapper pour les cinq à six mois qui viennent.
Cependant je ne peux le garantir, cela dépendra de Mr. Verdier, qui est à
Brandeis cette année. Je me mets en rapport avec lui ces jours-ci.

\page{3}
\note{double carbone d'une lettre de lui, dactylographiée}

Massy le 3.11.1970

Chère Madame Michon,

Je m'aperçois que dans les exemplaires SGA d'archives que j'ai emportés de chez
vous, il manque les suivants :

SGA 4 XVII, SGA 4 VI (première partie)

D'autre part j'ai besoin de ces exemplaires pour faire faire l'édition
photooffset en préparation. Pourriez vous me les retrouver ?

Bien cordialement

\page{4}
\note{double carbone d'une lettre de lui, dactylographiée}

Massy le 3.11.1970

Cher Deligne,

Pourrais-tu me dire si tu as terminé de regarder l'exposé de Rim SGA 7 VI, et si
oui, me l'envoyer avec tes commentaires (sinon, me dire si tu as l'intention de
le regarder et quand) ?

J'ai demandé à Mlle Altazin, qui a tapé ton exposé SGA 4 XVIII, de te l'envoyer
avec le manuscrit, pour que tu le corriges. Pourrais-tu me dire si tu l'as reçu
et si tu as l'intention de faire les corrections ? Quand elles seront faites, ou
si tu ne veux pas les faire, envoyes moi le texte au net stp.

As-tu eu des nouvelles de Ferrand pour SGA 6 XI ?

Bien cordialement

\section*{SGA 1}
\note{inscrit au crayon en haut à droite du feuillet de garde (page 5), qui ne
porte rien d'autre ; vraisemblablement de sa main. Le feuillet ne reçoit pas
de numéro de page}

\page{6}
\note{feuillet dactylographié, paginé « (ix) » à la machine : la dernière page
d'une préface de lui, dans un état antérieur au texte imprimé en tête de SGA 1
(Lecture Notes 224). La comparaison avec la préface publiée montre que ce n'est
pas le texte publié : celle-ci porte « immoral et dangereux » (avec une note
(4) dégageant la responsabilité de Springer), « quitte l'IHES le 30 septembre
1970 », « J'ai demandé à M. Motchane, directeur de l'IHES, que l'IHES s'abstienne
à partir du 1er octobre 1970 de diffuser \ldots », et ajoute un remerciement à
J.~P. Delale pour les index. Donné en entier ; rien de sa main sur le feuillet}

de ces questions que j'avais fait ailleurs. Ainsi, je me suis trouvé travailler
pendant trois ans dans une institution alors qu'elle participait à mon insu à un
mode de financement que je considère comme profondément immoral et extrêmement
dangereux. Etant à présent seul à avoir cette opinion parmi mes collègues à
l'IHES, ce qui a condamné à l'échec mes efforts pour obtenir la suppression des
subventions militaires du budget de l'IHES, j'ai pris la décision qui s'imposait
et quitte l'IHES le 30 Septembre prochain, et suspends également toute
collaboration scientifique avec cette Institution, aussi longtemps qu'elle
continuera à accepter de telles subventions. Monsieur L. Motchane, directeur de
l'IHES, s'est déclaré d'accord pour qu'à partir du 1er Octobre 1970, l'IHES ne
diffuse plus les textes mathématiques dont je suis auteur, ou faisant partie du
Séminaire de Géométrie Algébrique du Bois Marie. Comme il a été dit plus haut,
la diffusion de ce Séminaire va être assurée par la maison Julius Springer, dans
la série des Lecture Notes. Je suis heureux de remercier ici la maison Springer
et Monsieur K. Peters pour l'aide efficace et courtoise qu'ils m'ont apportée
pour rendre possible cette publication, en se chargeant en particulier de la
frappe pour la photooffset des nouveaux exposés rajoutés aux anciens séminaires,
et des exposés manquants des séminaires incomplets.

Massy, Août 1970

\section*{SGA 4}
\note{inscrit au crayon en haut à droite du feuillet de garde (page 7), qui ne
porte rien d'autre ; vraisemblablement de sa main. Le feuillet ne reçoit pas
de numéro de page}

\page{8}
\note{lettre dactylographiée de J.-L. Verdier (paraphe identique à celui de la
page 13), datée à la main « 16 Juillet », sans année, non reproduite. Il envoie
« la nouvelle rédaction de 4.10.6 » (pages 9 à 11), relève des « broutilles »
dans le n° 9 de SGA 4 IV (p. 149, 158, 159 du tirage) et dans les n° 10 et
suivants (p. 2 et 14), et signale qu'il manque le n° 14 (« Modules sur un topos
défini par recollement ») dans les notes qu'il a reçues, numéro qu'il se
souvient avoir rédigé. De sa main, au crayon, dans la marge gauche, en regard
de ce dernier paragraphe qu'un trait vertical signale : « demander à Messing »
(souligné)}

\page{9}
\note{première page (sur trois) du tapuscrit joint à la lettre de Verdier :
Exercice 4.10.6 de SGA 4 IV (un $\mathcal{U}$-site $S$ de topologie moins fine
que la canonique, une partie $\underline{M}$ de $\mathrm{Fl}(S)$ satisfaisant
a) à d), le « petit site » $S(X)$, questions 1) à 3)), dans la nouvelle
rédaction de Verdier ; non reproduit. Symboles complétés à l'encre ($\varphi_u$,
le soulignement de $\underline{M}$). Seule marque manuscrite : en marge gauche,
un signe d'insertion à l'encre devant « le site », tapé en marge, renvoyant à un
signe semblable après « ("petit site de X") » ; main non établie}

\page{10}
\note{deuxième page du même tapuscrit, non reproduite : fin de 3), les
foncteurs $\mathrm{Res}_X$ et $\mathrm{Prol}_X$, 4) le morphisme de topos
$f : S(X)^{\sim} \to (S_{/X})^{\sim}$, plongement du « petit topos » dans le
« gros topos », 5) et 6) sur la cohomologie. Marques manuscrites à l'encre : en
marge gauche, en regard de 5), « (V 4.1) » précédé d'un signe d'insertion qui
renvoie après « $S(X)$-acyclique » ; en fin de ligne, « $S(X)^{\sim}$ » écrit à
la main après « le petit topos ». Main non établie}

\page{11}
\note{troisième page du même tapuscrit, non reproduite : Exercice 4.10.7
(hypothèse (Et), le foncteur de changement de base $S(u)$, la rétraction
$g : S^{\sim}_{/X} \to S(X)^{\sim}$ du plongement $f$, le contre-exemple de la
topologie f.p.p.f., et un 4) « Acheter une médaille en chocolat pour les
rédacteurs »). Rien de sa main}

\page{12}
\note{photocopie d'une lettre de lui, rognée des deux côtés : les lettres
perdues aux bords sont restituées entre crochets}

Massy le 23.6.1971

\supplied{C}her Verdier,

Merci pour ta lettre du 23 Mai. C'est dommage que \supplied{tu}
\supplied{n}e m'aies pas envoyé ce malheureux exer 4.10.6, cela
a\supplied{urait} \supplied{p}ermis d'envoyer enfin à l'imprimeur le fascicule
1 de \supplied{S}GA 4. La personne qui a frappé le texte termine
mainte\supplied{nant}, \supplied{e}lle partira en vacances et Springer rend la
machine \supplied{q}u'elle avait en location, faute d'autres manuscrits.
\supplied{C}ela remet donc la publication même de ce fascicule sine
\supplied{d}ie.

Bien cordialement

\note{la lettre de Verdier à laquelle il répond est datée « le 28 Mai » (page
13) ; « 23 Mai » est bien ce que porte la page}

\page{13}
\note{lettre manuscrite de J.-L. Verdier sur papier à en-tête de Brandeis
University, Department of Mathematics, datée « le 28 Mai », sans année,
non reproduite. Elle répond à la lettre du 18.5.1971 (page 14) : l'appendice
« Image inverse d'un plat » ira à l'exposé V et l'appendice « Assez de points »
à l'exposé VI, ils ne concernent donc pas l'exposé IV ; il enverra « une version
non canulée de 4.10.6 » ; il demande une copie de IV §9--13 ; en marge, un
P.S. : il demandera à Mike de s'occuper des exposés après VI. Le reste est
personnel. De sa main, vraisemblablement : un ovale au crayon, très pâle,
autour de « de 4.10.6 »}

\page{14}
\note{double carbone d'une lettre de lui, dactylographiée ; la lettre se
poursuit sur un autre feuillet}

Massy le 18.5.1971

Cher Verdier,

Je ne trouve pas une minute pour m'occuper de la réédition de SGA 4. Saavedra
m'a un peu aidé, pour faire un index pour le Chap IV, par. 9 à 13. Mais on est
bloqué faute d'avoir les deux appendices de Deligne (assez, de points, image
inverse d'un plat), qu'il faudrait donner à la frappe. Les as-tu, et si oui,
peux-tu les envoyer pour la frappe ? D'autre part, peux-tu envoyer une version
corrigée de l'exercice 4.10.6 (médaille chocolat), j'ai entièrement oublié ce
qui était canulé. C'est un exercice qui sert tout le temps (!), il ne devrait
donc pas être question de le vider. Il doit y avoir quelque chose de vrai à
mettre à la place du faux \ldots

Si tu en as marre de t'occuper de la mise au point « définitive » de SGA 4, je
ne pourrais t'en blâmer. Je voudrais alors des instructions sur ce qu'il faut
faire -- je suis disposé à le laisser incomplet et dégueulasse, plûtôt que
d'avoir à y consacrer un jour ou deux aux dépens du boulot survie, infiniment
plus important. Pourrais-tu demander à Artin ce qu'il faut faire avec les
exposés après VI -- j'ai des tas de notes d'errata, footnotes améliorations
diverses -- d'en faire des instructions compréhensibles pour une dactylo
demanderait des jours de travail. Veut-il que je lui envoie mon paquet de
notes ?

\page{15}
\note{double carbone d'une lettre de lui, dactylographiée}

Massy le 3.11.1970

Cher Verdier,

Ne m'étant guère occupé de Math depuis trois mois, je suis un peu perdu pour
SGA 4. Si je me rappelle bien, tu as rédigé ou es en train de rédiger les
parties suivantes, qui sont exactement ce qui manque pour que SGA 4 soit
complet :

\struck{IV, par. 9 et suivants}\note{ligne biffée à la machine}

V

VI par. 5 et suivants

Si je me rappelle bien, \uncertain{VI}\note{chiffre surfrappé à la machine :
« V » ou « VI »} était en fait terminé d'être rédigé, il fallait seulement y
apporter quelques modifications dont on avait discuté avant ton départ. De
plus, je pense que tu as avec toi l'exposé

VI B

de Saint Donat, et je t'envoie également l'Appendice à XVII du même, dont
j'avais apparemment lu les premières pages. Pourrais tu me le renvoyer (ou le
renvoyer à St Donat) avec tes annotations et commentaires ?

Tu dois avoir un exemplaire du tirage de XVII ; le XVIII n'a pas été tapé sur
Stencils, mais directement sur papier pour offset ; la frappe est terminée, et
Deligne est censé la corriger.

Ecris-moi stp où tu en es avec ta part de rédaction, et avec la lecture de VI
B. Pour des raisons techniques, ce serait bien agréable si les textes V, VI,
$\mathrm{VI}_{\mathrm{B}}$ pouvaient être disponibles dans les deux, trois mois
qui viennent. Dis-moi en tous cas si tu as l'intention de terminer, et si oui,
quand tu penses avoir terminé.

Bien cordialement

\page{16}
\note{manuscrit de lui, à l'encre, sur papier à en-tête de l'Institut des
Hautes Études Scientifiques ; écriture rapide}

\emph{Continuation de} SGA 4 IV 4.10.6 (\add{généralisation de} 3° et 4°
appliqué à $S(X) \subset S_{/X}$)
\note{« généralisation de » est écrit au-dessus, une accolade le rattachant à
« 3° et 4° »}

5°) Soit $S_0 \hookrightarrow S$ une \struck{famille}
\struck{($S_0$ \ill{} \ill{})} \struck{et supposons que} sous-catégorie pleine
de $S$ satisfaisant les conditions suivantes :

(i) Si $u : X \to Y$ avec $u \in \underline{M}$, $Y \in \operatorname{ob} S_0$,
alors $X \in \operatorname{ob} S_0$.

[(ii) \struck{$\forall X \in \operatorname{ob} S$ existe famille couvrante
$X_i \to X$,} \struck{\ill{} \ill{}} le \uncertain{crible} de $S$ engendré par
les objets de $S_0$ est couvrant, i.e.\ $\forall X \in \operatorname{ob} S$,
existe famille couvrante $X_i \to X$, avec $X_i \to Z_i \in \operatorname{ob}
S_0$.]

Munissons $S_0$ de la top.\ induite. \uncertain{Prouvons} que \uncertain{une}
famille de morphismes $X_i \to X$ dans $S_0$ est couvrante dans $S_0$
\uncertain{ssi} elle l'est dans $S$ \struck{\ill{}} \uncertain{i.e.} que
$S_0 \to S$ est \emph{\uncertain{cocontinu}}.

Soit $\underline{M}_0$ l'ens.\ des $u \in \underline{M}$, \struck{\ill{}}
$u : X \to Y$, avec $Y \in \operatorname{ob} S_0$, donc
$X \in \operatorname{ob} S_0$. Prouvons que $\underline{M}_0 \subset
\mathrm{Fl}(S_0)$ satisfait les conditions a) b) c) d) de 1°).

Procédant comme dans 3°, définissons des foncteurs
\[
  \mathrm{Res} : S^{\sim} \to S_0^{\sim} ,
\]
\marginal{$\mathrm{Prol} : S_0^{\sim} \to S^{\sim}$ ? \uncertain{avec}
\uncertain{limites}}
\struck{et \uncertain{prouvons} que} \add{\uncertain{tels} que} Res
\uncertain{commute} \ill{} \uncertain{inductives} et \uncertain{projectives}, et
Prol \ill{} exact : \uncertain{p.\ f.} et \uncertain{pl.\ fidèle}, \ill{}
\ill{}
\note{la feuille s'arrête ici. Le « Prol » est dans un ovale tracé d'un trait
plus fin, à droite de « Res » ; la fin de la page (quatre lignes) est très
cursive et les lectures proposées sont guidées par le 3°) de l'exercice
(page 10 : $\mathrm{Res}_X$ commute aux limites inductives et projectives,
$\mathrm{Prol}_X$ pleinement fidèle)}

\page{17}
\note{feuillet dactylographié, que l'on donne en entier comme un tapuscrit de
lui, inédit (voir l'en-tête du fichier pour les raisons) : ce sommaire et sa
NB diffèrent du sommaire publié de l'exposé VI (qui a une introduction,
« 5. Commutation des foncteurs $H^i(X,-)$ aux limites inductives
filtrantes », et ni « Constructibilité » ni NB), et l'auteur y annonce qu'il
rédigera lui-même le n° 4, tandis que sa lettre du 3.11.1970 (page 15) confie à
Verdier « VI par. 5 et suivants ». Ses interventions à l'encre bleue sont entre
coins ; les surfrappes de la machine ne sont pas relevées}

\marginal{1.9.5, \uncertain{1.23}, 2.9 --- à reporter au \uncertain{n° 3}
\uncertain{3.1} ?}
\note{au crayon, très pâle, en tête du feuillet}

SGA 4 VI

\emph{Sommaire}

\begin{itemize}
\item[1.] Conditions de finitude pour les objets et flêches d'un topos. ---
\add{1 à 19}
\item[2.] Conditions de finitude pour un topos. --- \add{20 à 37}
\item[3.] Conditions de finitude pour un morphisme de topos.
\struck{\add{\ill{}}} --- \add{38 à 46}
\item[4.] \struck{Constructibilité.} \add{Conditions de finitude pour un topos
\add{obtenu par} recollement.} --- \add{47 à 65}
\item[\struck{5}\add{6}.] Catégories fibrées et $\varinjlim$ (rappels).
\item[\struck{6}\add{7}.] Sites et topos fibrés.
\item[\struck{7}\add{8}.] Limites projectives de topos.
\item[8.] Cas des limites projectives filtrantes de topos cohérents par
morphismes de transition cohérents : théorèmes de passage à la limite pour la
cohomologie.
\item[\struck{9}\add{5}.] Cohomologie d'une limite inductive de faisceaux sur
un topos cohérent.
\item[\add{10.}] \add{Appendice (par P. Deligne).} \struck{\add{Existence de
points \ill{}}} \add{topos cohérent\ill{}} \add{\uncertain{suffisamment}
\ill{} de points}
\end{itemize}
\note{les numéros de pages de la colonne de droite, à l'encre, sont coiffés
d'une abréviation soulignée illisible ; « 38 » est surchargé et relié par un
trait au n° 3. Le n° 4 dactylographié est barré à l'encre et remplacé
au-dessus, à la main ; un trait au crayon court sous la ligne. Les numéros 5,
6, 7 sont surchargés à l'encre en 6, 7, 8 ; le numéro du n° 8 est surchargé
d'un chiffre illisible, et une accolade au crayon réunit les n° 8 et 8 ; le 9
est surchargé en 5, entouré au crayon, et une longue flèche au crayon le fait
remonter après le n° 4. La ligne 10 est entièrement manuscrite, marquée d'un
trait vertical au crayon en marge ; après « Appendice (par P. Deligne). » un
début de titre est biffé et récrit, avec une ligne interlinéaire, d'une
écriture très rapide}

\struck{NB Bien sûr 8 serait interchangeable avec le bloc \uncertain{5} à
\uncertain{7}. Mais il me semble plus naturel de mettre dans deux numéros
consécutifs les deux facettes du passage à la limite en cohomologie. Bien sûr,
le numéro 8 sera réduit à une ou deux pages, le sorite des conditions de
finitude étant déjà fait.}

\struck{Le n° 4 contiendra le sorite des groupes constructibles (au sens
groupes finis ou non), des faisceaux de torsion, des notions de
constructibilité pour les Modules et complexes de Modules, -- y compris les
faits spéciaux dans le cas localement noethérien. Je pense que je le rédigerai
mais ne le garantis pas. On peut laisser tomber le n°4, ce sera alors moins
beau bien sûr \ldots :}
\note{ces deux paragraphes sont annulés par des traits obliques croisés au
crayon ; dans le premier, les chiffres « 8 » et « 5 à 7 » sont surchargés à
l'encre (lecture des chiffres d'origine douteuse)}

La rédaction qui suit des n°s 1 à 3 est complète, à cela près que je compte y
ajouter encore quelques exemples (en particulier un topos localement algébrique
non algébrique\add{)}, \struck{un morphisme non cohérent de topos
cohérents).} \add{et un topos ayant une famille génératrice formée d'objets
cohérents, qui n'est pas localement algébrique}
\note{l'ajout, au crayon, occupe les deux dernières lignes du feuillet ; un
double trait au crayon le signale dans la marge gauche}

\page{18}
\note{notes manuscrites, à l'encre noire puis bleue}

SGA 4 \emph{XVII}\note{le numéro est surligné et souligné}

\begin{itemize}
\item[p. 75 l. $-4$] réf à \emph{V} correcte ?
\item[p. 90 l. 2] réf : $f_{!}$ pour Ensembles pointés resp.\ faisceaux en
groupes ?? (Consulter Verdier)
\item[p. 139 l. $-10$] réf : $\mathrm{VI}_{\mathrm{B}}$
\item[p. 167 l. $-8$] réf : EGA II 2e édition (lemme de
Chow)
\item[170 l. 7] réf : \emph{V} \ill{} Deligne \ldots
\item[175 l. $-1$] id
\end{itemize}
\note{une accolade réunit les lignes 170 et 175, dont la référence est
commune}

\page{19}
\note{notes manuscrites, à l'encre noire et bleue ; plan de SGA 4 en trois
fascicules}

\emph{SGA 4}

\emph{Fasc.} 1 : Préface, avant-propos, table des matières générale du Sém.
Exp I à \emph{IV}. Table des matières du Volume, index
\uncertain{terminologique} \uncertain{général}. $\#\,570$ p.\note{le 5 de « 570 » est repris à l'encre ; la ligne « Exp. V VI VII
VIII, IX X XI » d'en dessous, à droite, est reliée par un trait au total}

\emph{Fasc.} 2 : \uncertain{Topologie générale} ; \uncertain{Cohomologie des
topos et des sites étalés} : \uncertain{étude générale}. Exp.\ V, VI
($\mathrm{VI}_{\mathrm{B}}$), VII, VIII, IX, X, XI.
\note{les titres du fascicule 2 sont écrits en oblique, soulignés, d'une
écriture très rapide ; lectures douteuses}

\begin{itemize}
\item[V] 100
\item[VI] 150
\item[$\mathrm{VI}_{\mathrm{B}}$] 30
\item[VII] 24
\item[VIII] 46
\item[IX] \struck{\ill{}} 42
\item[X] 21
\item[XI] 19
\end{itemize}
Total : 432\note{« 100 » récrit à l'encre noire sur un premier chiffre}

\emph{Fasc.} 3 : \struck{Cat.\ \uncertain{Dérivées} \ill{}} Théorèmes de
\uncertain{changement} de base et de dualité globale\note{une ligne partant de
sous « $\mathrm{VI}_{\mathrm{B}}$ 30 » et une autre partant de sous « XI 19 »
rejoignent ce titre ; à la seconde est attaché, biffé, « \ill{} », et, dans un
ovale, « Cat.\ dér.\ ? 40 »}

\begin{itemize}
\item[XII] 57
\item[XIII] 15
\item[XIV] 24
\item[XV] 38
\item[XVI] 44
\item[XVII] 210
\item[XVIII] 200
\item[XIX] 51
\end{itemize}
$\#\,640$
\note{une accolade réunit XII à XVI, avec au crayon « 178 » (leur somme) ; un
« 5 » isolé au crayon plus bas}

\page{20}
\note{notes manuscrites, à l'encre bleue, avec des additions au crayon}

\struck{Il faut \ill{}} \emph{insérer le texte manquant en}
\struck{\emph{une}} p. 83 bis, dans laquelle il faudra reporter le texte des
lignes 7 à 18 du \struck{\ill{}} \add{stencil} 84

\uncertain{me montrer} la frappe faite

\begin{itemize}
\item[p. 85 l. 13] $\mathbb{L}g^{*} \quad \mathbb{L}g'^{*}$
\item[l. 23] $\mathbb{L}g'^{*} \quad \mathbb{L}g^{*}$
\item[l. 39] $\mathbb{L}g'^{*}$
\item[l. 27 et suite :]
\end{itemize}
\note{il écrit les deux formes côte à côte, sans flèche : la première à
remplacer par la seconde, comme on le comprend ; « $\mathbb{L}$ » rend son L à
double hampe}

\struck{Le texte oublié a été reporté ici \emph{par erreur} !!!}
\note{les lignes « Il faut insérer \ldots » jusqu'à « \ldots par erreur !!! »
sont barrées de trois longs traits obliques au crayon ; la dernière phrase,
soulignée, est à l'encre sous les traits}

$210 + 76 + 250 = 536$
\note{addition posée en colonne, au crayon, le total entouré ; au-dessus de 76, un premier
chiffre (4 ?) est barré, et la somme vaut 536 avec 76. Un « Fasc. » isolé
à droite}

\end{document}
